problem:  
Let \((M^n,g,e^{-\phi} dV_g)\) be a complete smooth metric–measure space satisfying the Bakry–Émery curvature–dimension condition  
\[
\mathrm{Ric}_{\phi,N} := \mathrm{Ric}_g + \nabla^2 \phi - \frac{1}{N-n}\nabla \phi\otimes \nabla \phi \ge (N-1)K g
\]
for some \(N>n\) and \(K>0\).  
Let \(\Sigma\) be a smooth, closed, embedded hypersurface bounding a domain \(\Omega\), and define the weighted mean curvature \(H_{\phi} = H - \langle\nabla \phi ,\nu\rangle.\)

It is known that various Heintze–Karcher-type inequalities hold under curvature–dimension constraints, but only for specific model densities (e.g. Gaussian or exponential).  

**Conjecture / Problem:**  
Is there a **dimension–sharp and curvature–sharp Heintze–Karcher inequality** of the form  
\[
\int_{\Sigma}\frac{e^{-\phi}}{H_{\phi}}\,d\mu \geq \Psi_{K,N}\big(\mathrm{Vol}_{\phi}(\Omega)\big)
\]
where \(\Psi_{K,N}\) is an explicit monotone function depending only on \(K,N\), such that  
1. equality holds **only** on the \(N\)-dimensional weighted model space of constant curvature \(K\) (for instance, the spherical, Gaussian, or exponential space form), and  
2. the inequality remains **valid in the synthetic curvature–dimension setting** \(CD(K,N)\) for general metric measure spaces?  

Equivalently, can one establish such a sharp inequality purely under the \(CD(K,N)\) (possibly \(RCD(K,N)\)) conditions without smoothness assumptions, thereby determining whether Heintze–Karcher–type comparison principles characterize the curvature–dimension structure in a metric–measure sense?

Why is it a "good" problem:  
This revised problem formulates a precise next frontier beyond current literature. Weighted Heintze–Karcher inequalities for smooth Bakry–Émery manifolds exist in partial or nonsharp forms, but a **dimension–sharp, curvature–sharp** version valid for all \(CD(K,N)\) spaces (including nonsmooth ones) is still open. Such a result would unify classical Riemannian comparison theorems, weighted Reilly-type formulas, and synthetic geometry under a single inequality principle, bridging **geometric analysis, optimal transport**, and **metric geometry**.  
Its conceptual simplicity (a single inequality) hides deep analytical difficulties—it would likely require new variational and transport tools to generalize curvature comparisons to the nonsmooth, weighted context. Success would extend the foundational reach of the Heintze–Karcher inequality into the heart of the Lott–Sturm–Villani curvature–dimension theory, exemplifying the hallmark of a “good” research problem: simple to state, profound in potential impact, and currently unsolved.